\chapter{The raw view}\label{ch:raw}
\begin{goals}
\textbf{In this chapter} you will look at any file's bytes as they are, as lines or as hex, and move
between a row in the grid and the bytes it was read from.
\end{goals}

\begin{pending}
This chapter waits on the raw view, which is being built (the workspace plan,
\file{spec\_inbox/lamina-workspace-spec.md}, §3). It will cover:
\begin{itemize}
\item \textbf{Lines}: line numbers and byte offsets; CR, LF, CRLF, tabs, the BOM and control characters
made visible; bytes that don't decode in the chosen encoding marked; no wrapping, horizontal scroll.
\item \textbf{Hex}: sixteen bytes a row, offset · hex · text, for any file, binary ones included.
\item \textbf{Search}: the find bar in both; in lines, the filter as grep, and export of what it keeps.
\item \textbf{The jumps}: from a grid row to its bytes (\ui{Show raw}), from a line back to its row,
between lines and hex at the same offset.
\item \textbf{Compressed and in-memory sources}: offsets in decompressed space; a workbook sheet's
container bytes.
\item \textbf{The interpretation dialog's raw pane}: setting the delimiter, header and skipped lines
while looking at the real first lines.
\item \textbf{Data quality into raw}: from a flag to the bytes it is about.
\end{itemize}
\end{pending}
